\section{Fill time as a survival problem}
\label{sec:survival}

A fill rule answers \emph{whether} an order would be filled given the market that
followed. A strategy deciding where to quote needs the answer \emph{before} the
market moves: how soon an order placed now is likely to be filled. Most orders
are withdrawn before they could have been filled (Section~\ref{sec:outcomes}), so
fill time is a time-to-event outcome under heavy censoring, and the factorisation
\eqref{eq:factor} splits it into two such outcomes.

\subsection{Design}

\paragraph{Stages.}
\emph{Front}: from placement to the front of the queue, for every order.
\emph{Fill}: from the front to the first fill, for orders that reached the front.
\emph{Total}: from placement to the first fill. Each is censored by the owner's
cancellation or by the end of observation; hours are rebuilt between their
neighbours, so an order placed late in an hour is followed into the next.

\paragraph{Features.}
Only what was known just before the order arrived: the book one change earlier
and trades strictly before placement. Side; order size and volume ahead (logs);
whether the level was empty; distance behind the best price on the order's side
and from the mid-price, in ticks; spread; queue imbalance at the best prices
\citep{gould2016queue}, signed towards the order's side; trade blocks that hit the
order's side in the last 1 and 10~seconds; and the mid-price move towards the
order over the same windows.

\paragraph{Model and evaluation.}
A Cox proportional-hazards model per stage, ridge-penalised ($0.01$), is fitted on
the orders of hours 0--11 and scored on hours 12--23 of the same day by Harrell's
concordance $C$---the share of comparable pairs of orders whose predicted order of
events is the observed one. One million orders are drawn at random from each half.
As a benchmark we take the ranking of the exponential clock: $\lambda(\delta)$
depends on the distance from the mid alone, so its concordance is that of $\delta$.
The three event rules are not comparable in this way: they read the market after
placement, which a predictive model cannot.

\subsection{Results}

\begin{table}[t]
\centering\small
\caption{Concordance on held-out hours (12--23) of a Cox model fitted on hours
0--11, and of single features: the volume ahead, and the distance from the mid,
which is the exponential clock's ranking.}
\label{tab:survival}
\setlength{\tabcolsep}{4pt}
\begin{tabular}{@{}ll ccc@{}}
\toprule
 & Stage & Cox & Volume ahead & Clock ($\delta$) \\
\midrule
\inputrows{tables/survival_rows.tex}
\bottomrule
\end{tabular}
\end{table}

\begin{table}[t]
\centering\small
\caption{Hazard ratios of the Cox model for the total time to fill, fitted on
hours 0--11. A ratio above one shortens the expected time to fill.}
\label{tab:hazard}
\setlength{\tabcolsep}{4pt}
\begin{tabular}{@{}l ccc@{}}
\toprule
Feature & \BTC & \ETH & \SOL \\
\midrule
\inputrows{tables/hazard_rows.tex}
\bottomrule
\end{tabular}
\end{table}

\paragraph{Prediction.}
On hours it never saw, the model ranks the total time to fill with a concordance
of 0.858 on \SOL{} and 0.813 on \ETH{} (Table~\ref{tab:survival}), against 0.775
and 0.773 for the exponential clock and 0.807 and 0.766 for the volume ahead
alone. Concordance on the test hours is 0.02--0.03 below that on the training
hours, which come earlier in the same day. On \BTC: \tbd.

\paragraph{What shortens the wait.}
Every unit of log volume ahead lowers the fill hazard by 17--19\%, and joining an
empty level raises it by a factor of 1.25 (\SOL) and 2.48 (\ETH)
(Table~\ref{tab:hazard}). Larger orders wait longer. Recent trading on the
order's side raises the hazard by 5--6\% per trade block in the last second and
2--3\% per block in the last ten. Each tick of spread raises it, and each tick of
distance from the best price lowers it slightly. Imbalance towards the order's
own side lowers it, consistent with queue imbalance predicting the next move away
from the heavier side \citep{gould2016queue}. Side does not matter.

\paragraph{The two stages.}
The factors of \eqref{eq:factor} have different drivers. Reaching the front is
a matter of the queue: the volume ahead alone ranks it with $C=0.853$ (\SOL) and
0.811 (\ETH), while the distance from the mid is close to uninformative (0.53--0.54).
Once an order leads its level, the distance from the mid takes over (0.61 and
0.67), and the volume that was ahead at placement ranks the remaining wait
\emph{backwards} ($C=0.45$ and 0.46): orders that waited behind long queues are
filled sooner once they lead. A plausible reading is that a long queue is consumed
mostly when the price is moving to the level, so an order that reaches the front
of one meets an aggressor almost at once. A single intensity per distance, as in the exponential clock,
cannot represent either stage.
